\hypertarget{project-objective-and-current-progress}{%
\section{Project objective and current
progress}\label{project-objective-and-current-progress}}

My goal is to implement the computational methods in \emph{Pooling
Strategy for Massive Viral Sequencing} and evaluate sample recovery
using public HCV HVR1 sequences. The problem is to design overlapping
pools that identify individual specimens, then assign the sequences
observed in those pools back to their source specimens.

I have implemented the compatibility graph, Greedy Pool Design Algorithm
(GPDA), tabu-search clique-cut selection, kGEM clustering, and
generalized pool intersections and differences. I have also completed an
exploratory evaluation using eight public specimens. The implementation
is a chapter-based reproduction with documented interpretations of the
printed pseudocode, rather than a claim of exact reproduction of the
published experiments.

\hypertarget{methodology-public-hcv-dataset}{%
\section{Methodology: public HCV
dataset}\label{methodology-public-hcv-dataset}}

\hypertarget{data-source-and-specimen-selection}{%
\subsection{Data source and specimen
selection}\label{data-source-and-specimen-selection}}

I used the HCV HVR1 sequence deposit associated with Janiak et
al.~(2019), available on Zenodo under DOI 10.5281/zenodo.3266795 and the
CC BY 4.0 license. The source contains reconstructed, error-corrected
haplotypes and estimated frequencies in FASTA headers. These are not raw
sequencing reads.

Before evaluating recovery, I selected the first eight files labeled
chronic genotype 1 in numeric specimen-ID order: \textbf{53, 80, 142,
153, 157, 165, 200, and 201}. The rule is deterministic but does not
produce a random or representative cohort. The filenames use genotype 1;
the study describes its genotype-1 cohort as subtype 1b. I did not
independently re-genotype the sequences.

\begin{longtable}[]{@{}lr@{}}
\toprule\noalign{}
Public specimen & Source haplotype records \\
\midrule\noalign{}
\endhead
\bottomrule\noalign{}
\endlastfoot
53 & 16 \\
80 & 70 \\
142 & 35 \\
153 & 24 \\
157 & 20 \\
165 & 5 \\
200 & 54 \\
201 & 114 \\
\textbf{Total} & \textbf{338} \\
\end{longtable}

\hypertarget{preparation-and-pooling-protocol}{%
\subsection{Preparation and pooling
protocol}\label{preparation-and-pooling-protocol}}

I retained all listed haplotypes without an additional frequency cutoff
and normalized the header frequency estimates within each specimen. I
removed the source gap symbols and aligned all sequences together using
MAFFT 7.526 with \texttt{-\/-auto\ -\/-thread\ 1}. The alignment
contains 180 columns and 333 distinct strings. Checksums verify the
downloaded files, and sequence-content checks confirm that alignment did
not change the ungapped nucleotide strings.

For this evaluation, I assumed all eight specimens were mutually
compatible. Measured titers and epidemiological exclusions were not
used. Thus, the experiment tests deconvolution under a complete
compatibility graph, rather than validating the biological compatibility
model.

I used a maximum pool size of four and independently sampled 80
observations per included specimen per pool using the normalized source
frequencies. Each replicate has 1,040 observations. I ran five
replicates using seeds 91 through 95. The initial clustering multiplier
was W=2; deconvolution was repeated from the original pools with
increased W, up to W=4, if any recovered specimen was empty.

\textbf{The input haplotypes are real public HCV sequences, but pooling
and read sampling are simulated.} No laboratory pooling, new sequencing,
or added sequencing-error model was used. Ground-truth source labels
were withheld from clustering and used only to evaluate assignments.
Independent observations retain distinct identifiers, even if their
nucleotide strings are identical.

\hypertarget{algorithm}{%
\section{Algorithm}\label{algorithm}}

\hypertarget{pool-design-with-gpda-and-tabu-search}{%
\subsection{Pool design with GPDA and tabu
search}\label{pool-design-with-gpda-and-tabu-search}}

Let G represent sample compatibility and H represent unresolved sample
pairs. Initially, H contains every pair of specimens. A candidate pool A
is valid if it is a clique in G and its size does not exceed the limit
T=4. It separates a pair when exactly one endpoint belongs to A.

The chapter's clique-cut objective can be written as

\[
f(A)=c_H(A,V\setminus A)-(|E(H)|+1)q_G(A),
\]

where c\_H counts unresolved pairs separated by A and q\_G counts
incompatible pairs inside A. The penalty discourages invalid pools. The
tabu heuristic searches both orientations of a MaxCut initialization,
uses the chapter's single-vertex gain formulas to make improving moves,
and uses undo operations and temporarily forbidden moves to explore
alternatives. It trims an oversized record to satisfy T. The
implementation also uses documented iteration limits and a singleton
fallback.

GPDA adds the selected pool and removes the pairs it separates from H.
It repeats until H has no edges, then adds an uncovered singleton if
needed. Finally, it checks that every specimen is covered and has a
distinct pool-membership signature.

For the public specimens, the resulting design was:

\begin{longtable}[]{@{}ll@{}}
\toprule\noalign{}
Pool & Specimens included \\
\midrule\noalign{}
\endhead
\bottomrule\noalign{}
\endlastfoot
P1 & 53, 142, 157, 200 \\
P2 & 53, 80, 157, 165 \\
P3 & 53, 80, 142, 153 \\
P4 & 201 \\
\end{longtable}

This gives four pools for eight specimens. The same design is used
across the five sampling replicates. GPDA and tabu search are
heuristics; I do not claim that their output is globally optimal. The
additional exact clique-cut solver is used only for small-instance
verification.

\hypertarget{kgem-clustering}{%
\subsection{kGEM clustering}\label{kgem-clustering}}

For two pools, I set the requested cluster count to k=W times the number
of specimens in their membership union. If k exceeds the number of
distinct observed strings, the implementation caps it. The initial
centroid is the most frequent read, and subsequent seeds maximize their
minimum Hamming distance to the existing seeds.

For centroid g\_i and aligned read r of length L, the emission
likelihood is

\[
h_{i,r}=\prod_{m=1}^{L}f_{i,m}(r_m).
\]

The consensus allele at a position has probability 1-4 epsilon and each
of the other four alleles has probability epsilon; the alphabet includes
a deletion symbol. The current implementation uses epsilon=0.01. Given
mixture proportions pi\_i, posterior cluster membership is

\[
p_{i,r}=\frac{\pi_i h_{i,r}}{\sum_j\pi_j h_{j,r}}.
\]

The inner iteration estimates mixture proportions using observed read
counts. The outer iteration updates the consensus alleles using the
chapter's printed posterior sum over distinct read strings. This is a
discrete consensus update, not a full soft allele-frequency update.
Likelihoods are computed in logarithmic form for numerical stability.
The iterations stop when estimates stabilize or a documented iteration
or cycle limit is reached. Each observation is assigned to its
highest-posterior cluster.

\hypertarget{generalized-intersections-and-deconvolution}{%
\subsection{Generalized intersections and
deconvolution}\label{generalized-intersections-and-deconvolution}}

For pools P and Q, I cluster their combined observations and retain
clusters that contain reads from both pools. These clusters form the
generalized intersection. The known pool memberships determine which
specimens this intersection represents; the nucleotide identities
determine which reads enter it. Directed differences remove the selected
observation IDs from the corresponding original pool.

The queue procedure repeatedly replaces a pair of overlapping pools by
its intersection and nonempty directed differences until individual
specimen sets remain. Replacing both operands is an explicit
interpretation of the chapter's queue pseudocode. Increasing W is
triggered only by an empty recovered specimen; a nonempty but incorrect
result does not trigger another attempt.

\hypertarget{results-on-the-public-hcv-dataset}{%
\section{Results on the public HCV
dataset}\label{results-on-the-public-hcv-dataset}}

Accuracy is the number of correct assignment entries divided by all
assignment entries. Read recall is the number of correctly recovered
observation IDs divided by all input observation IDs. Every observation
was assigned exactly once in these runs, so accuracy and recall
coincide.

\begin{longtable}[]{@{}lr@{}}
\toprule\noalign{}
Measure & Public HCV result \\
\midrule\noalign{}
\endhead
\bottomrule\noalign{}
\endlastfoot
Specimens & 8 \\
Source haplotype records & 338 \\
Alignment columns & 180 \\
Pools & 4 \\
Observations per replicate & 1,040 \\
Mean assignment accuracy and read recall & 71.81\% \\
Accuracy range across five seeds & 70.48--72.88\% \\
Mean per-specimen recall (macro average) & 72.97\% \\
Mean recall excluding the singleton specimen & 69.46\% \\
\end{longtable}

\hypertarget{replicate-results}{%
\subsection{Replicate results}\label{replicate-results}}

\begin{longtable}[]{@{}
  >{\raggedright\arraybackslash}p{(\columnwidth - 6\tabcolsep) * \real{0.2000}}
  >{\raggedleft\arraybackslash}p{(\columnwidth - 6\tabcolsep) * \real{0.2667}}
  >{\raggedleft\arraybackslash}p{(\columnwidth - 6\tabcolsep) * \real{0.2667}}
  >{\raggedleft\arraybackslash}p{(\columnwidth - 6\tabcolsep) * \real{0.2667}}@{}}
\toprule\noalign{}
\begin{minipage}[b]{\linewidth}\raggedright
Sampling seed
\end{minipage} & \begin{minipage}[b]{\linewidth}\raggedleft
Correct assignments / observations
\end{minipage} & \begin{minipage}[b]{\linewidth}\raggedleft
Accuracy
\end{minipage} & \begin{minipage}[b]{\linewidth}\raggedleft
Final W
\end{minipage} \\
\midrule\noalign{}
\endhead
\bottomrule\noalign{}
\endlastfoot
91 & 746 / 1,040 & 71.73\% & 3 \\
92 & 758 / 1,040 & 72.88\% & 3 \\
93 & 747 / 1,040 & 71.83\% & 4 \\
94 & 750 / 1,040 & 72.12\% & 4 \\
95 & 733 / 1,040 & 70.48\% & 3 \\
\end{longtable}

The variation here comes from sampling the same underlying specimens
repeatedly. It does not measure uncertainty across independent cohorts.
All 48 kGEM fits in the main runs, including retries, reported inner and
outer convergence. This shows numerical stabilization, not correctness
of the resulting specimen assignments.

\hypertarget{specimen-level-recovery}{%
\subsection{Specimen-level recovery}\label{specimen-level-recovery}}

For seed 91, recovery differed substantially between specimens:

\begin{longtable}[]{@{}lrr@{}}
\toprule\noalign{}
Public specimen & Correct / true observations & Recall \\
\midrule\noalign{}
\endhead
\bottomrule\noalign{}
\endlastfoot
53 & 240 / 240 & 100.00\% \\
80 & 6 / 160 & 3.75\% \\
142 & 70 / 160 & 43.75\% \\
153 & 80 / 80 & 100.00\% \\
157 & 160 / 160 & 100.00\% \\
165 & 50 / 80 & 62.50\% \\
200 & 60 / 80 & 75.00\% \\
201 & 80 / 80 & 100.00\% \\
\end{longtable}

Specimen 201 occupies a singleton pool, so its complete recovery does
not test separation of mixed populations. Excluding that specimen gives
a more informative additional measure of performance on pooled
specimens. Specimen 80 remains poorly recovered despite a nonempty
output, showing why the empty-sample retry rule is not sufficient to
detect failure.

\hypertarget{sensitivity-to-w}{%
\subsection{Sensitivity to W}\label{sensitivity-to-w}}

I also used seed 91 with retries disabled to examine fixed W values:

\begin{longtable}[]{@{}
  >{\raggedright\arraybackslash}p{(\columnwidth - 6\tabcolsep) * \real{0.2143}}
  >{\raggedleft\arraybackslash}p{(\columnwidth - 6\tabcolsep) * \real{0.2857}}
  >{\raggedleft\arraybackslash}p{(\columnwidth - 6\tabcolsep) * \real{0.2857}}
  >{\raggedright\arraybackslash}p{(\columnwidth - 6\tabcolsep) * \real{0.2143}}@{}}
\toprule\noalign{}
\begin{minipage}[b]{\linewidth}\raggedright
Fixed W
\end{minipage} & \begin{minipage}[b]{\linewidth}\raggedleft
Correct / total observations
\end{minipage} & \begin{minipage}[b]{\linewidth}\raggedleft
Accuracy
\end{minipage} & \begin{minipage}[b]{\linewidth}\raggedright
Empty recovered specimens
\end{minipage} \\
\midrule\noalign{}
\endhead
\bottomrule\noalign{}
\endlastfoot
1 & 499 / 1,040 & 47.98\% & 80, 157 \\
2 & 710 / 1,040 & 68.27\% & 80 \\
3 & 746 / 1,040 & 71.73\% & None \\
\end{longtable}

The results depend on W, but these three runs do not establish that
increasing W always improves recovery. They are a sensitivity analysis
on one sampling replicate, not a separate validation set or a
parameter-optimization study.

\hypertarget{verification-and-limitations}{%
\section{Verification and
limitations}\label{verification-and-limitations}}

The saved full test run contains 19 passing tests. Relevant checks
include pool feasibility and separation, clique-cut gain formulas,
source-file checksums, frequency parsing, preservation of ungapped
sequence content, and metric accounting under deliberately incorrect
assignments.

The main limitations of the present evaluation are:

\begin{itemize}
\tightlist
\item
  The dataset contains reconstructed haplotypes and frequency estimates,
  not raw reads or experimentally pooled specimens. The simulation does
  not include PCR bias, contamination, unequal titers, or sequencing
  errors.
\item
  Only eight chronic genotype-1 specimens from one study were evaluated.
  All samples were assumed compatible, so the experiment does not test
  measured biological compatibility restrictions.
\item
  One aligned sequence occurs in more than one specimen. Identical
  strings alone cannot uniquely establish the original source of an
  observation.
\item
  Numerical convergence and nonempty recovered sets do not guarantee
  correct biological assignments. Poor recovery can remain undetected by
  the retry condition.
\item
  The implementation includes documented repairs to stopping conditions,
  bounded search, and queue refinement. It is not a literal execution of
  every printed pseudocode line or a reproduction of the chapter's
  published empirical accuracy.
\item
  Four pools instead of eight individually processed specimens
  demonstrates a reduction in pool count in this setup. It does not
  establish a 50\% monetary saving, because total sequencing depth and
  laboratory costs were not held fixed.
\end{itemize}

\hypertarget{next-steps}{%
\section{Next steps}\label{next-steps}}

My next priority is to trace the generalized intersections that
misassign observations from specimens 80 and 142. I plan to inspect the
saved confusion matrices and clustering diagnostics, then vary W,
sampling depth, and pool-size limits separately. I also plan to evaluate
additional public specimens using selection rules set before inspecting
their results. Any alternative consensus update or queue strategy will
be identified as a departure from the current chapter-based
implementation. These are planned investigations, not completed
improvements.

\hypertarget{references-and-supporting-results}{%
\section{References and supporting
results}\label{references-and-supporting-results}}

\begin{enumerate}
\def\labelenumi{\arabic{enumi}.}
\tightlist
\item
  Skums et al.~(2016). \emph{Pooling Strategy for Massive Viral
  Sequencing}. Chapter 3 in \emph{Computational Methods for Next
  Generation Sequencing Data Analysis}, pp.~57--83.
\item
  Janiak et al.~(2019). \emph{Hepatitis C virus (HCV) genotype 1b
  displays higher genetic variability of hypervariable region 1 (HVR1)
  than genotype 3}.
  \href{https://doi.org/10.1038/s41598-019-49258-y}{Scientific Reports}.
\item
  Caraballo Cortes, K. (2019). Public HCV HVR1 sequence dataset.
  \href{https://doi.org/10.5281/zenodo.3266795}{Zenodo, DOI
  10.5281/zenodo.3266795}, CC BY 4.0.
\end{enumerate}

\begin{itemize}
\tightlist
\item
  \href{https://github.com/trilokpadhi/approx-algorithms-viral-pooling/blob/main/data/public_hcv/README.md}{Public-data
  provenance and preparation}
\item
  \href{https://github.com/trilokpadhi/approx-algorithms-viral-pooling/blob/main/results/public_hcv/report-results.json}{Public
  HCV results and diagnostics}
\item
  \href{https://github.com/trilokpadhi/approx-algorithms-viral-pooling/blob/main/results/verification.txt}{Saved
  full test output}
\end{itemize}
